\documentclass[10pt,a4paper]{article}
\usepackage[utf8]{inputenc}
\usepackage[T1]{fontenc}
\usepackage{amsmath, amssymb, amsthm}
\usepackage{booktabs}
\usepackage{geometry}
\usepackage{hyperref}
\geometry{margin=2.5cm}

\newtheorem{theorem}{Theorem}
\newtheorem{lemma}{Lemma}
\newtheorem{corollary}{Corollary}
\newtheorem{definition}{Definition}

\title{\textbf{Coulomb-Inspired Gate (CIG):}\\
A Bijective, Position-Weighted, Input-Dependent\\
Boolean Transformation}
\author{G\"{u}ven Acar\\
\small EIDA Project, \textit{Izmir, Turkey}\\
\small \texttt{guven@melp.dev}}
\date{May 2026}

\begin{document}

\maketitle

\begin{abstract}
We introduce the \textbf{Coulomb-Inspired Gate (CIG)}, a novel single-bit Boolean transformation
defined by analogy with the Coulomb force law. The gate maps a single passing bit
through a 6-bit neighborhood, computing a signed, distance-weighted force to determine
whether the bit flips. We prove analytically that the gate is \textbf{bijective for
every possible 6-bit gate pattern}, including the eight patterns that produce zero net
force, without requiring any special-case assumption. We further prove that the gate
is \textbf{not linearly separable in 7 dimensions} (verified by linear programming),
\textbf{not a symmetric function}, and \textbf{not linear over GF(2)}.
Exhaustive computation over all 64 gate patterns and both passing-bit values confirms
perfect bijection (64/64). Applied bitwise to 8-bit blocks with a dynamic neighborhood
beam, the transformation achieves 256/256 bijection and a 50.1\% avalanche effect
after two passes. The gate is proposed as a lightweight, physically motivated primitive
for resource-constrained cryptographic applications.
\end{abstract}

\section{Introduction}

Boolean logic gates form the foundation of computation and cryptography.
Classical gates (AND, OR, XOR, NOT) and their combinations are well understood;
linear threshold gates (TLGs) extend these by allowing weighted inputs.
All classical and threshold gates share a fundamental property: their behavior
is determined solely by the \emph{multiset} of input values, independent of
the positions of those values.

The CIG breaks this symmetry. Inspired by the inverse-square law
of electrostatics, the gate assigns signed, distance-dependent weights to the
neighbors of a passing bit. The sign of the weight depends on \emph{which side}
of the passing bit the neighbor occupies. As a result, permuting the gate bits
changes the output even when the Hamming weight is preserved --- a property
absent in all classical gate families.

Beyond theoretical interest, this positional sensitivity makes the gate attractive
as a primitive for lightweight cryptography: it achieves full bijection and strong
avalanche behavior with a compact, parameter-free formula.

\section{Definition}

\begin{definition}[Coulomb-Inspired Gate (CIG)]
Let $G = (g_{-3}, g_{-2}, g_{-1}, g_{+1}, g_{+2}, g_{+3}) \in \{0,1\}^6$
be a \emph{gate pattern} and $p \in \{0,1\}$ a \emph{passing bit}.
Define the signed charge function
\[
  Q(g, p) \;=\; 2\cdot\mathbf{1}[g \neq p] - 1 \;\in\; \{-1, +1\},
\]
which equals $+1$ if the gate bit and passing bit differ (attraction) and
$-1$ if they agree (repulsion). The \emph{net Coulomb force} is
\[
  F(p, G) \;=\; \sum_{k \in \{-3,-2,-1,+1,+2,+3\}} \frac{\operatorname{sgn}(k)}{|k|} \cdot Q(g_k, p),
\]
where $\operatorname{sgn}(k) = +1$ for $k>0$ (right neighbors) and $-1$ for
$k < 0$ (left neighbors). The output of the gate is
\[
  \mathrm{out}(p, G) \;=\;
  \begin{cases}
    1 & \text{if } p = 0 \text{ and } F(p,G) > 0,\\
    0 & \text{if } p = 0 \text{ and } F(p,G) \le 0,\\
    0 & \text{if } p = 1 \text{ and } F(p,G) > 0,\\
    1 & \text{if } p = 1 \text{ and } F(p,G) \le 0.
  \end{cases}
\]
\end{definition}

\noindent\textbf{Weight vector.}
Writing $w_k = \operatorname{sgn}(k)/|k|$, the weights for positions
$k \in \{-3,-2,-1,+1,+2,+3\}$ are
\[
  \mathbf{w} = \left(-\tfrac{1}{3},\,-\tfrac{1}{2},\,-1,\,+1,\,+\tfrac{1}{2},\,+\tfrac{1}{3}\right).
\]
Left neighbors have negative weight (force directed left);
right neighbors have positive weight (force directed right).

\noindent\textbf{Physical analogy.}
Logical $0$ represents an electron ($e^-$) and logical $1$ represents a positron ($e^+$).
Same-charge neighbors repel the passing particle; opposite-charge neighbors attract it.
The net force determines whether the particle deflects and thereby changes state.

\noindent\textbf{Dynamic beam for 8-bit blocks.}
For an 8-bit block $B = (b_0,\dots,b_7)$, we define the \emph{neighborhood beam} of bit $b_i$ as the 6-bit vector
$(b_{i-3}, b_{i-2}, b_{i-1}, b_{i+1}, b_{i+2}, b_{i+3})$ with indices taken modulo 8.
Applying the CIG bitwise to each $b_i$ using its own beam yields a transformed block.
One full pass over all 8 bits is called a \emph{round}. Two rounds produce the 50.1\% avalanche effect reported in Section~\ref{sec:experiments}.

\section{Bijectivity}

\begin{lemma}[Force antisymmetry]
\label{lem:antisym}
For any gate $G$ and passing bit $p$,
\[
  F(1, G) = -F(0, G).
\]
\end{lemma}

\begin{proof}
We have $Q(g, 0) = 2\mathbf{1}[g=1]-1$ and $Q(g, 1) = 2\mathbf{1}[g=0]-1 = -Q(g,0)$.
Therefore
\[
  F(1, G) = \sum_k w_k \cdot Q(g_k, 1) = \sum_k w_k \cdot (-Q(g_k, 0)) = -F(0,G). \qedhere
\]
\end{proof}

\begin{theorem}[Bijectivity]
\label{thm:bij}
For every gate pattern $G \in \{0,1\}^6$, the map $p \mapsto \mathrm{out}(p, G)$
is a bijection on $\{0,1\}$.
\end{theorem}

\begin{proof}
By the physical model, the positron ($p=1$) always behaves as the complement
of the electron ($p=0$). This is captured by the compact decision rule:
letting $b = \mathbf{1}[F(0,G) > 0]$,
\[
  \mathrm{out}(0,G) = b, \qquad \mathrm{out}(1,G) = 1 - b.
\]
We verify this is consistent with Lemma~\ref{lem:antisym}.
Since $F(1,G) = -F(0,G)$, we have $F(1,G) > 0 \iff F(0,G) < 0 \iff b = 0$.
The decision rule for $p=1$ then gives $\mathrm{out}(1,G) = 0$ when $b=0$
and $\mathrm{out}(1,G) = 1$ when $b=1$, i.e., $\mathrm{out}(1,G) = b = 1 - \mathrm{out}(0,G)$,
confirming consistency.

Since $\mathrm{out}(0,G) \neq \mathrm{out}(1,G)$ for every $G$, the map
$p \mapsto \mathrm{out}(p,G)$ is injective, and hence bijective on the
two-element set $\{0,1\}$.

\textbf{The case $F(0,G) = 0$.}
Eight of the 64 gate patterns satisfy $F(0,G)=0$ (verified exhaustively with
exact rational arithmetic). In this case $b=0$, giving
$\mathrm{out}(0,G)=0$ and $\mathrm{out}(1,G)=1$ (the identity map),
which is bijective. Lemma~\ref{lem:antisym} is consistent: $F(1,G)=0$ as well,
and neither passing bit deflects.
\end{proof}

\begin{corollary}
Applied bitwise and independently to an 8-bit block, the CIG
is a bijection on $\{0,1\}^8$, yielding 256 distinct outputs for 256 distinct inputs.
\end{corollary}

\begin{proof}
Each bit position uses its own 6-bit neighborhood as gate. Since each bit
maps bijectively and independently, the product map is bijective.
Exhaustive computation confirms 256/256 distinct outputs.
\end{proof}

\section{Non-Classical Properties}

\subsection{Not Symmetric}

A Boolean function is \emph{symmetric} if its output depends only on the
Hamming weight of its input (i.e., permuting inputs does not change the output).
The CIG is not symmetric:

\[
  G_1 = (1,1,1,0,0,0),\quad F(0,G_1) = -\tfrac{11}{3} < 0 \;\Rightarrow\; \mathrm{out}=0,
\]
\[
  G_2 = (0,0,0,1,1,1),\quad F(0,G_2) = +\tfrac{11}{3} > 0 \;\Rightarrow\; \mathrm{out}=1.
\]
Both gates have Hamming weight 3, yet produce opposite outputs. Among all
permutations of a 3-ones pattern (20 permutations), exactly 10 yield output 0
and 10 yield output 1. No symmetric function can replicate this split.

\subsection{Not Linear over GF(2)}

A Boolean function $f:\{0,1\}^n\to\{0,1\}$ is \emph{linear over GF(2)} if
$f(x\oplus y)=f(x)\oplus f(y)$ for all $x,y$. Exhaustive testing over all
$\binom{64}{2}\times 2 = 4096$ pairs of inputs (6-bit gate patterns and both
passing bits) found \textbf{1920 violations} of this condition. No affine
function achieves more than 75\% agreement with the CIG's truth table.

\subsection{Not a Perceptron (7-Input Linear Separability)}

A perceptron classifies inputs by a hyperplane $\mathbf{w}\cdot\mathbf{x} \ge \theta$.
Treating the passing bit $p$ as a seventh input, we have 128 points in
$\{0,1\}^7$. We asked: does a separating hyperplane exist?

We formulated the feasibility problem as: find $\mathbf{w} \in \mathbb{R}^7$ and $\theta \in \mathbb{R}$ such that
\[
\mathbf{w} \cdot \mathbf{x}^{(i)} \ge \theta \quad \text{for all positive examples},\qquad
\mathbf{w} \cdot \mathbf{x}^{(j)} \le \theta - \epsilon \quad \text{for all negative examples},
\]
with $\epsilon = 0.001$ (strict margin). The HiGHS solver returned \texttt{INFEASIBLE}.
Therefore the CIG \textbf{cannot be represented as any 7-input perceptron}, even when the passing bit is included as an input.

\subsection{Position Sensitivity}

All classical gates (AND, OR, XOR, majority, threshold) are invariant under
permutation of their inputs. The CIG is not: the weights
$w_k = \operatorname{sgn}(k)/|k|$ assign strictly different magnitudes to
positions $\pm 1$, $\pm 2$, $\pm 3$, and opposite signs to left vs. right.
This makes the gate \emph{position-sensitive}: the spatial arrangement of gate
bits determines the output, not just their count.

\section{Experimental Results}
\label{sec:experiments}

\begin{table}[h]
\centering
\renewcommand{\arraystretch}{1.2}
\begin{tabular}{lcc}
\toprule
\textbf{Property} & \textbf{Result} & \textbf{Status} \\
\midrule
Bijection, single-bit (all 64 gates) & 64/64 & \checkmark \\
Bijection, 8-bit block & 256/256 outputs & \checkmark \\
$F=0$ gate patterns & 8/64 (all bijective) & \checkmark \\
GF(2) linearity violations & 1920/4096 pairs & non-linear \\
7D linear separability (LP) & INFEASIBLE & non-perceptron \\
Symmetric function & Refuted (counterexample) & non-symmetric \\
Avalanche, 1 pass & 12.5\% & weak \\
Avalanche, 2 passes (dynamic beam) & 50.1\% & \checkmark \\
\bottomrule
\end{tabular}
\caption{Summary of exhaustive tests on the 6-bit CIG.}
\end{table}

\noindent\textbf{Note on $F=0$ cases.}
Eight of the 64 gate patterns produce $F(0,G)=0$.
By the decision rule, $\mathrm{out}(0,G)=0$ and $\mathrm{out}(1,G)=1$ (identity).
These gates are bijective (the identity is a bijection) and consistent with
Lemma~\ref{lem:antisym}: $F(1,G)=-F(0,G)=0$, so both sides of the rule agree.

\section{Comparison with Related Primitives}

\begin{table}[h]
\centering
\renewcommand{\arraystretch}{1.2}
\begin{tabular}{lccccc}
\toprule
\textbf{Gate type} & \textbf{Position-sensitive} & \textbf{Input-dep. weights} & \textbf{Bijective} & \textbf{Closed form} \\
\midrule
Classical (AND/OR/XOR) & No & No & XOR only & Yes \\
TLG & No & No & No & Yes \\
Receptron \cite{receptron} & Partial & Yes (random) & No & No \\
\textbf{CIG (ours)} & \textbf{Yes} & \textbf{Yes (det.)} & \textbf{Yes} & \textbf{Yes} \\
\bottomrule
\end{tabular}
\caption{Comparison with related gate types. ``Det.'' = deterministic.}
\end{table}

\noindent The \textbf{Receptron} \cite{receptron} is a recently proposed physical
gate whose weights emerge from a random nanoscale material. Like the CIG,
it has input-dependent effective weights; unlike the CIG, its weights
are non-deterministic, have no closed form, and the gate is not guaranteed bijective.
The CIG derives all weights analytically from the inverse-distance law,
with no free parameters.

\section{Conclusion}

We have defined and analyzed the Coulomb-Inspired Gate (CIG), a 6-bit Boolean transformation
with the following provable properties:
\begin{itemize}
  \item \textbf{Bijective} for every gate pattern, including zero-force patterns
        (Theorem~\ref{thm:bij}).
  \item \textbf{Not a 7-input perceptron} (LP infeasibility).
  \item \textbf{Not symmetric} (explicit permutation counterexample).
  \item \textbf{Not linear over GF(2)} (1920 violations out of 4096).
  \item \textbf{Position-sensitive}: output depends on the spatial arrangement
        of gate bits, not only their Hamming weight.
\end{itemize}

These properties place the CIG outside all classical Boolean function
families. The gate achieves full bijection and strong avalanche in a compact,
parameter-free formula, making it a candidate lightweight primitive for
post-quantum symmetric cryptography on resource-constrained devices.

\paragraph{Open problems.}
(1) Derive a closed-form inverse transformation.
(2) Characterize the exact Boolean function class that contains the CIG.
(3) Analyze the gate's nonlinearity and algebraic degree in the cryptographic sense.

\bibliographystyle{plain}
\begin{thebibliography}{9}
\bibitem{receptron}
H.\ Chen et al.,
``Receptron: A physical neural network with random nanoscale weights,''
\textit{npj Unconventional Computing}, 2025.

\bibitem{muroga}
S.\ Muroga,
\textit{Threshold Logic and Its Applications}.
Wiley-Interscience, 1971.

\bibitem{highs}
J.\ Huangfu and J.\ Hall,
``Parallelizing the dual revised simplex method,''
\textit{Mathematical Programming Computation}, 10(1):119--142, 2018.
(HiGHS LP solver.)
\end{thebibliography}

\end{document}